% !TeX root = ../papers/nc-p3.tex
\section{二直線列の両端中央支持の全制限}
\label{sec:two-line-full-middle-restriction}
\begin{theorem}\label{thm:two-line-full-middle-restriction}
曲面射影のない標準正則窓の二直線列を $p,L,C,z$ とする。
$p$ の中央支持の $C$ への制限と $z$ の中央支持の $L$ への制限はいずれも階数2である。
$L,C$ 上の全点の中央支持は4である。
両端の中央支持次元を $a,b$、$r=\operatorname{rank}w(p,z)$ とすると
\[
             1\le r\le\min(2,a-2,b-2),\qquad \rho\ge6+a+b-r .
\]
従って二直線列には $\rho\ge11$ が必要で、等号なら $a=b=3$, $r=1$ である。
特に反対階数9,10の分岐では純関係を許しても全ての列が少なくとも三つの動く位置を持つ。
\end{theorem}
\begin{proof}
$L,C$ の消去二空間を $S,M$ とする。$P=R_1\cap(S\otimes V_2)$ と
$Q=R_1\cap(V_1\otimes M)$ はそれぞれ三次元で、$w(p),w(z)$ の末尾、先頭の縮約像である。
$w(p)$ の中央支持を $D\subset V_2$ とする。$R_2$ の $C$ への制限は
命題\ref{prop:two-line-branch-block}の束の列によって単射なので $D$ の $C$ への制限は非零である。
さらに $w(p)$ を $C$ に制限しても、位置1の $S$ 支持は二次元のままである。
実際、位置1での縮約の二空間が $R_2$ に入り、この単射で二次元を保つ。

もし $D$ の制限が一次元なら、制限した $w(p)$ はその固定中央文字と両端の行列の積で、
この行列は左支持が二次元なので階数2である。従って $P$ の $C$ への制限の階数は2で、
$P\cap Q$ の次元は高々1。ところが $D$ の制限を消す $y\in C$ があり、
$P+Q$ の少なくとも五関係が $y$ で消えてしまう。
六次元の $R_1$ の逆評価が階数1以下になり、定理\ref{thm:all-endpoint-ranks-two}に反する。
これで制限階数2が従い、逆向きにも同じである。

任意の $x\in L$ に対し、$w(-,x,C,-)$ は $[x,y]h$ の表示により
位置3で $\ker z$ の全三空間を張る。さらに $w(x,z)\ne0$ である。
これは $z$ の中央支持の $L$ への制限の階数が2だからである。
従って $w(x)$ の中央支持は4であり、逆側も同様である。
$w(p,z)$ の位置2支持は $D\cap M$ に含まれその上界は $a-2$ となる。
位置1側の上界も同じである。$w(p,z)=0$ なら $z\to p$ も辺になり、異なる
パラメータの $p,L,C,z$ で二つの四閉路を自由に反復でき、標準成長に反する。
あとは命題\ref{prop:two-line-branch-block}の切断不等式から従う。
\end{proof}
